%%%%%%%%%%%%%%%%%%%%%%%%%%%%%%%%%%%%%%%%%%%%%%%%%%%%%%%%%%%%%%%%%%%%%%%%%%%%%%%%%%
\begin{frame}[fragile]\frametitle{}

\begin{center}
{\Large AIOOLL Workshop}
\end{center}
\end{frame}

%%%%%%%%%%%%%%%%%%%%%%%%%%%%%%%%%%%%%%%%%%%%%%%%%%%%%%%%%%%%%%%%%%%%%%%%%%%%%%%%%%
\begin{frame}[fragile]\frametitle{Workshop Agenda}
\begin{columns}
\column{0.48\textwidth}
\textbf{Hour 1: Setup \& Foundations}
\begin{itemize}
  \item[\textbf{00:00}] What is AIOOLL?
  \item[\textbf{00:10}] Install Ubuntu on old laptop
  \item[\textbf{00:25}] System prep \& Miniconda
  \item[\textbf{00:35}] Ollama install \& model pull
  \item[\textbf{00:45}] Clone repo \& bootstrap
  \item[\textbf{00:55}] Module 1: Classical ML
\end{itemize}
\column{0.48\textwidth}
\textbf{Hour 2: AI Modules}
\begin{itemize}
  \item[\textbf{01:10}] Module 2: Deep Learning (PyTorch)
  \item[\textbf{01:25}] Module 3: RAG Pipeline
  \item[\textbf{01:40}] Module 4: AI Agents
  \item[\textbf{01:50}] Module 5: Computer Vision
  \item[\textbf{01:55}] Module 6: LLM Inference
  \item[\textbf{02:00}] Tests \& Next Steps
\end{itemize}
\end{columns}
\vspace{0.4em}
\begin{block}{Hardware You Need}
Intel i3 laptop (any generation) $\cdot$ 8\,GB RAM $\cdot$ SSD (strongly preferred) $\cdot$ USB drive (8\,GB+)
\end{block}
\end{frame}

%%%%%%%%%%%%%%%%%%%%%%%%%%%%%%%%%%%%%%%%%%%%%%%%%%%%%%%%%%%%%%%%%%%%%%%%%%%%%%%%%%
\begin{frame}[fragile]\frametitle{}
\begin{center}
{\Large \textbf{Part 0: What is AIOOLL?}}

\textit{Understanding the philosophy before touching the keyboard}
\end{center}
\end{frame}

%%%%%%%%%%%%%%%%%%%%%%%%%%%%%%%%%%%%%%%%%%%%%%%%%%%%%%%%%%%%%%%%%%%%%%%%%%%%%%%%%%
\begin{frame}[fragile]\frametitle{What is AIOOLL?}
\begin{itemize}
  \item \textbf{AIOOLL} = \textbf{A}I \textbf{O}n \textbf{O}ld \textbf{L}inux \textbf{L}aptop
  \item \textbf{The Claim:} A complete AI development stack: classical ML, deep learning, RAG, agents, computer vision, and local LLMs: running entirely on a CPU-only laptop with 8\,GB RAM.
  \item \textbf{The Philosophy:} AI should be accessible to everyone. Cloud credits and gaming GPUs are not prerequisites for real AI work.
  \item \textbf{Key Technology:} \textit{Model quantization} (GGUF format, Q4\_K\_M) reduces a 2B-parameter LLM from 4\,GB to 1.5\,GB while retaining most capability. The \texttt{ollama} runtime handles all of this automatically.
  \item \textbf{What You Will Build Today:} Six working AI applications: each with a Python driver, pytest test suite, and a Streamlit UI: all offline, all on your laptop.
\end{itemize}
\begin{block}{Real Benchmark: Intel i3-8130U, 8\,GB RAM, SSD}
  Gemma2:2b chat: \textbf{$\sim$8 tok/s} $\cdot$ RAG query: \textbf{$\sim$50\,ms} $\cdot$ ML train: \textbf{$<$1\,s} $\cdot$ CV face detect: \textbf{15 FPS}
\end{block}
\end{frame}

%%%%%%%%%%%%%%%%%%%%%%%%%%%%%%%%%%%%%%%%%%%%%%%%%%%%%%%%%%%%%%%%%%%%%%%%%%%%%%%%%%
\begin{frame}[fragile]\frametitle{AIOOLL: Six Modules at a Glance}
\begin{center}
\begin{tabular}{|p{1.5cm}|p{1.5cm}|p{2cm}|p{1.8cm}|}
\toprule
\textbf{Module} & \textbf{Tech} & \textbf{Task} & \textbf{Time} \\
\midrule
1. Classical ML & scikit-learn & Spam detect, price predict & $<$1\,s \\
2. Deep Learning & PyTorch (CPU) & LSTM sentiment, MLP regression & $\sim$30\,s/ epoch \\
3. RAG Pipeline & LangChain + ChromaDB & Document Q\&A, local KB & $\sim$50\,ms/ query \\
4. AI Agents & LangGraph + Ollama & Research \& code review agents & $\sim$5--15\,s \\
5. Computer Vision & OpenCV & Face detect, motion detect & 15\,FPS \\
6. LLM Inference & Ollama & Benchmark, chat, prompting & $\sim$8\,tok/s \\
\bottomrule
\end{tabular}
\end{center}
\vspace{0.3em}
\textbf{Every module provides:}
\begin{multicols}{2}
\begin{itemize}
  \item \texttt{driver.py}: headless runner
  \item \texttt{tests/test\_*.py}: pytest suite
  \item \texttt{ui/app.py}: Streamlit dashboard
  \item Bundled dataset (no download)
\end{itemize}
\end{multicols}
\end{frame}

%%%%%%%%%%%%%%%%%%%%%%%%%%%%%%%%%%%%%%%%%%%%%%%%%%%%%%%%%%%%%%%%%%%%%%%%%%%%%%%%%%
\begin{frame}[fragile]\frametitle{}
\begin{center}
{\Large \textbf{Part 1: Install Ubuntu}}

\textit{Prerequisite: Free your hardware from Windows}
\end{center}
\end{frame}

%%%%%%%%%%%%%%%%%%%%%%%%%%%%%%%%%%%%%%%%%%%%%%%%%%%%%%%%%%%%%%%%%%%%%%%%%%%%%%%%%%
\begin{frame}[fragile]\frametitle{Why Xubuntu? }

Choosing the Right Flavour

\begin{itemize}
  \item \textbf{Standard Ubuntu (GNOME):} $\sim$800\,MB idle RAM: leaves only 7.2\,GB for AI workloads on an 8\,GB machine. \textbf{Avoid.}
  \item \textbf{Xubuntu (XFCE):} $\sim$280\,MB idle RAM: recommended. Lightweight, stable, full hardware support.
  \item \textbf{Lubuntu (LXQt):} $\sim$200\,MB idle: slightly lighter, less polished. Good alternative.
  \item \textbf{Version:} Ubuntu 22.04 LTS (Jammy Jellyfish): supported until April 2027. Stable Python 3.10 available natively.
  \item \textbf{Download Xubuntu ISO:} \href{https://xubuntu.org/download/}{\texttt{xubuntu.org/download}} (choose 22.04 LTS, 64-bit)
\end{itemize}
\begin{block}{RAM Savings Summary}
GNOME 800\,MB $\rightarrow$ XFCE 280\,MB = \textbf{520\,MB freed}: roughly the size of a quantized 500M-parameter model
\end{block}
\end{frame}

%%%%%%%%%%%%%%%%%%%%%%%%%%%%%%%%%%%%%%%%%%%%%%%%%%%%%%%%%%%%%%%%%%%%%%%%%%%%%%%%%%
\begin{frame}[fragile]\frametitle{Step 1: Create a Bootable USB Drive}
\begin{itemize}
  \item \textbf{Tool needed:} Rufus (Windows): free, at \href{https://rufus.ie}{\texttt{rufus.ie}}
  \item \textbf{USB drive:} 8\,GB minimum, all data will be erased
\end{itemize}
\vspace{0.3em}
\textbf{Procedure in Rufus:}
\begin{enumerate}
  \item Plug in USB drive
  \item Open Rufus $\rightarrow$ \textit{Device}: select your USB drive
  \item \textit{Boot selection}: click \textbf{SELECT} $\rightarrow$ choose downloaded \texttt{xubuntu-22.04-desktop-amd64.iso}
  \item \textit{Partition scheme}: \textbf{GPT} (modern laptops) or \textbf{MBR} (very old BIOS)
  \item \textit{File system}: FAT32
  \item Click \textbf{START} $\rightarrow$ accept ISO mode warning $\rightarrow$ wait $\sim$5 minutes
  \item Click \textbf{CLOSE} when done
\end{enumerate}
\begin{alertblock}{Important}
Back up all Windows data before proceeding. The next steps will modify your disk.
\end{alertblock}
\end{frame}

%%%%%%%%%%%%%%%%%%%%%%%%%%%%%%%%%%%%%%%%%%%%%%%%%%%%%%%%%%%%%%%%%%%%%%%%%%%%%%%%%%
\begin{frame}[fragile]\frametitle{Step 2: Boot from USB \& Install Xubuntu}
\textbf{Enter BIOS/UEFI boot menu:}
\begin{itemize}
  \item Restart laptop $\rightarrow$ press \textbf{F12} (Dell/Lenovo/HP) or \textbf{F2/Esc/Del} (varies by brand)
  \item Select USB drive from boot menu
\end{itemize}
\vspace{0.3em}
\textbf{Xubuntu Installer: key choices:}
\begin{enumerate}
  \item \textbf{Try or Install:} Choose \textit{Install Xubuntu}
  \item \textbf{Keyboard layout:} Select your language
  \item \textbf{Updates \& other software:} Tick ``Install third-party software'' (codecs, drivers)
  \item \textbf{Installation type:}
  \begin{itemize}
    \item \textit{Erase disk and install Xubuntu}: full replacement (simplest, recommended for this workshop)
    \item \textit{Install alongside Windows}: dual boot (keep Windows)
  \end{itemize}
  \item \textbf{Timezone, name, password:} Fill in and note your password
  \item Wait $\sim$10--20 minutes $\rightarrow$ \textbf{Restart Now} when prompted $\rightarrow$ remove USB
\end{enumerate}
\begin{block}{After Reboot}
Login with your password. You now have Xubuntu 22.04. Open a terminal: \textbf{Ctrl+Alt+T}
\end{block}
\end{frame}

%%%%%%%%%%%%%%%%%%%%%%%%%%%%%%%%%%%%%%%%%%%%%%%%%%%%%%%%%%%%%%%%%%%%%%%%%%%%%%%%%%
\begin{frame}[fragile]\frametitle{Step 3: First Boot: Update \& Essential Tools}
\begin{lstlisting}
# Update package lists and upgrade all packages
sudo apt update && sudo apt upgrade -y

# Install essential build tools
sudo apt install -y build-essential curl wget git \
    htop btop lm-sensors python3-pip

# Check sensors (CPU temperature monitoring)
sudo sensors-detect --auto
sensors

# Set CPU to performance mode (run before AI workloads)
echo performance | sudo tee \
    /sys/devices/system/cpu/cpu*/cpufreq/scaling_governor

# Verify Python version (should be 3.10)
python3 --version
\end{lstlisting}
\end{frame}

%%%%%%%%%%%%%%%%%%%%%%%%%%%%%%%%%%%%%%%%%%%%%%%%%%%%%%%%%%%%%%%%%%%%%%%%%%%%%%%%%%
\begin{frame}[fragile]\frametitle{Step 4: Install Miniconda}
\begin{lstlisting}
# Download Miniconda installer (Linux 64-bit)
wget https://repo.anaconda.com/miniconda/\
Miniconda3-latest-Linux-x86_64.sh \
    -O ~/miniconda.sh

# Run installer (-b = batch, -p = prefix path)
bash ~/miniconda.sh -b -p $HOME/miniconda3

# Initialise conda for bash
~/miniconda3/bin/conda init bash

# Reload shell (or open new terminal)
source ~/.bashrc

# Verify installation
conda --version          # should show: conda 24.x.x
conda info | head -5
\end{lstlisting}
\begin{block}{Why Conda?}
Conda manages Python versions per environment, isolates packages, and handles binary dependencies (numpy, torch) more reliably than pip alone on Linux.
\end{block}
\end{frame}

%%%%%%%%%%%%%%%%%%%%%%%%%%%%%%%%%%%%%%%%%%%%%%%%%%%%%%%%%%%%%%%%%%%%%%%%%%%%%%%%%%
\begin{frame}[fragile]\frametitle{}
\begin{center}
{\Large \textbf{Part 2: AIOOLL Setup}}\

\textit{Clone the repo, create the environment, pull models}
\end{center}
\end{frame}

%%%%%%%%%%%%%%%%%%%%%%%%%%%%%%%%%%%%%%%%%%%%%%%%%%%%%%%%%%%%%%%%%%%%%%%%%%%%%%%%%%
\begin{frame}[fragile]\frametitle{Clone Repo \& Create Conda Environment}
\begin{lstlisting}
# Clone the AIOOLL repository
git clone https://github.com/yourusername/AIOOLL.git
cd AIOOLL

# Create conda environment with Python 3.10,
# activate it, and install all dependencies
conda create -n aiooll python=3.10 -y \
    && conda activate aiooll \
    && pip install -r requirements.txt

# Run the full bootstrap script
# (downloads NLTK data + OpenCV Haar cascades)
bash scripts/setup.sh

# Verify the environment
conda activate aiooll
python -c "import sklearn, torch, cv2, streamlit; \
    print('All core packages OK')"
python -c "import torch; \
    print(f'PyTorch {torch.__version__}, CPU only')"
\end{lstlisting}
\end{frame}

%%%%%%%%%%%%%%%%%%%%%%%%%%%%%%%%%%%%%%%%%%%%%%%%%%%%%%%%%%%%%%%%%%%%%%%%%%%%%%%%%%
\begin{frame}[fragile]\frametitle{Install Ollama \& Pull Quantized Models}
\begin{lstlisting}
# Install Ollama (official installer)
curl -fsSL https://ollama.ai/install.sh | sh

# Start Ollama server in background
ollama serve &

# Pull quantized CPU-friendly models
# Gemma2 2B  ~1.5 GB  -- best for chat and RAG
ollama pull gemma2:2b

# Qwen2 1.5B  ~0.9 GB  -- fastest, multilingual
ollama pull qwen2:1.5b

# Embedding model for RAG  ~270 MB
ollama pull nomic-embed-text

# Verify models are ready
ollama list
\end{lstlisting}
\begin{alertblock}{Disk Space}
All three models together use $\approx$2.7\,GB. Ensure at least 5\,GB free on your SSD before pulling.
\end{alertblock}
\end{frame}

%%%%%%%%%%%%%%%%%%%%%%%%%%%%%%%%%%%%%%%%%%%%%%%%%%%%%%%%%%%%%%%%%%%%%%%%%%%%%%%%%%
\begin{frame}[fragile]\frametitle{Repository Structure: Quick Map}
\begin{lstlisting}[language={}]
AIOOLL/
 src/
  ml/            # Module 1: scikit-learn
  |-- driver.py  # run with: python driver.py
  |-- ui/app.py  # run with: streamlit run ui/app.py
  |-- tests/     # run with: pytest tests/ -v
  |-- data/      # sms_spam.csv, house_prices.csv
  dl/            # Module 2: PyTorch
  rag/           # Module 3: LangChain + ChromaDB
  agents/        # Module 4: LangGraph
  cv/            # Module 5: OpenCV
  llm_inference/ # Module 6: Ollama benchmarking
 tests/          # Integration tests (run from root)
 scripts/
  setup.sh       # bootstrap (run once)
  pull_models.sh # pull all Ollama models
 Makefile        # make run-ml / make ui-rag / etc.
\end{lstlisting}
\end{frame}

%%%%%%%%%%%%%%%%%%%%%%%%%%%%%%%%%%%%%%%%%%%%%%%%%%%%%%%%%%%%%%%%%%%%%%%%%%%%%%%%%%
\begin{frame}[fragile]\frametitle{Makefile: Your Shortcut Reference}
\begin{lstlisting}
# Always activate the environment first
conda activate aiooll

make setup        # install deps + NLTK + CV cascades
make models       # pull all Ollama models

# Run headless drivers
make run-ml       # cd src/ml && python driver.py
make run-dl       # cd src/dl && python driver.py
make run-rag      # cd src/rag && python driver.py
make run-agents   # cd src/agents && python driver.py
make run-cv       # cd src/cv  && python driver.py
make run-llm      # cd src/llm_inference && python driver.py

# Launch Streamlit UIs (opens browser automatically)
make ui-ml        # streamlit run src/ml/ui/app.py
make ui-rag       # streamlit run src/rag/ui/app.py
make ui-agents    # streamlit run src/agents/ui/app.py

make test         # pytest all modules
make clean        # remove all generated artifacts
\end{lstlisting}
\end{frame}

%%%%%%%%%%%%%%%%%%%%%%%%%%%%%%%%%%%%%%%%%%%%%%%%%%%%%%%%%%%%%%%%%%%%%%%%%%%%%%%%%%
\begin{frame}[fragile]\frametitle{}
\begin{center}
{\Large \textbf{Module 1: Classical Machine Learning}}\

\texttt{src/ml/driver.py} $\cdot$ scikit-learn $\cdot$ CPU $<$1\,s
\end{center}
\end{frame}

%%%%%%%%%%%%%%%%%%%%%%%%%%%%%%%%%%%%%%%%%%%%%%%%%%%%%%%%%%%%%%%%%%%%%%%%%%%%%%%%%%
\begin{frame}[fragile]\frametitle{Module 1: ML Concepts}
\begin{itemize}
  \item \textbf{Task A: Spam Detection:} Binary text classification using TF-IDF features fed into four algorithms: Naive Bayes, Logistic Regression, Linear SVM, Random Forest. Algorithms compete; best AUC-ROC wins.
  \item \textbf{Task B: House Price Prediction:} Tabular regression on 8 engineered features (area, age, distance, amenities). Ridge Regression vs.\ Gradient Boosting compared on $R^2$, RMSE, MAE.
  \item \textbf{TF-IDF:} Term Frequency × Inverse Document Frequency. Converts raw text to a numerical matrix where common words are down-weighted. \texttt{ngram\_range=(1,2)} captures both single words and two-word phrases.
  \item \textbf{Pipeline:} scikit-learn \texttt{Pipeline} chains vectorizer + classifier so the same transformations are applied consistently to train and test splits: preventing data leakage.
  \item \textbf{Cross-Validation:} 5-fold stratified CV gives an unbiased estimate of generalisation performance by training/testing on 5 different splits.
\end{itemize}
\begin{block}{Run It}
  \texttt{cd src/ml \&\& python driver.py} $\quad$ then $\quad$ \texttt{make ui-ml}
\end{block}
\end{frame}

%%%%%%%%%%%%%%%%%%%%%%%%%%%%%%%%%%%%%%%%%%%%%%%%%%%%%%%%%%%%%%%%%%%%%%%%%%%%%%%%%%
\begin{frame}[fragile]\frametitle{Module 1: ML Code: Imports \& Paths}
\begin{lstlisting}
import numpy as np
import pandas as pd
import matplotlib.pyplot as plt
from pathlib import Path
from sklearn.model_selection import (
    train_test_split, cross_val_score, StratifiedKFold
)
from sklearn.pipeline import Pipeline
from sklearn.preprocessing import StandardScaler, LabelEncoder
from sklearn.feature_extraction.text import TfidfVectorizer
from sklearn.metrics import (
    classification_report, roc_auc_score,
    mean_squared_error, r2_score, mean_absolute_error
)
from sklearn.naive_bayes import MultinomialNB
from sklearn.linear_model import LogisticRegression, Ridge
from sklearn.svm import LinearSVC
from sklearn.ensemble import (
    RandomForestClassifier, GradientBoostingRegressor
)
import joblib
from loguru import logger

ROOT     = Path(__file__).parent
DATA_DIR = ROOT / "data"
MODELS_DIR = ROOT / "models"; MODELS_DIR.mkdir(exist_ok=True)
RESULTS_DIR= ROOT / "results"; RESULTS_DIR.mkdir(exist_ok=True)
\end{lstlisting}
\end{frame}

%%%%%%%%%%%%%%%%%%%%%%%%%%%%%%%%%%%%%%%%%%%%%%%%%%%%%%%%%%%%%%%%%%%%%%%%%%%%%%%%%%
\begin{frame}[fragile]\frametitle{Module 1: ML Code: SpamClassifier (Part 1)}
\begin{lstlisting}
class SpamClassifier:
    ALGORITHMS = {
        "Naive Bayes": Pipeline([
            ("tfidf", TfidfVectorizer(
                ngram_range=(1, 2), max_features=5000,
                sublinear_tf=True)),
            ("clf", MultinomialNB(alpha=0.1))
        ]),
        "Logistic Regression": Pipeline([
            ("tfidf", TfidfVectorizer(
                ngram_range=(1, 2), max_features=5000,
                sublinear_tf=True)),
            ("clf", LogisticRegression(
                C=1.0, max_iter=1000, random_state=42))
        ]),
        "Linear SVM": Pipeline([
            ("tfidf", TfidfVectorizer(
                ngram_range=(1, 2), max_features=5000,
                sublinear_tf=True)),
            ("clf", LinearSVC(
                C=1.0, max_iter=1000, random_state=42))
        ]),
    }
    def __init__(self, data_path=DATA_DIR/"sms_spam.csv"):
        self.data_path = data_path
        self.results: dict = {}
        self.best_model = None
\end{lstlisting}
\end{frame}

%%%%%%%%%%%%%%%%%%%%%%%%%%%%%%%%%%%%%%%%%%%%%%%%%%%%%%%%%%%%%%%%%%%%%%%%%%%%%%%%%%
\begin{frame}[fragile]\frametitle{Module 1: ML Code: SpamClassifier (Part 2)}
\begin{lstlisting}
    def load_data(self):
        df = pd.read_csv(
            self.data_path, encoding="utf-8", quotechar='"')
        df["label_num"] = (df["label"] == "spam").astype(int)
        return df

    def run(self):
        df = self.load_data()
        X, y = df["text"], df["label_num"]
        X_train, X_test, y_train, y_test = train_test_split(
            X, y, test_size=0.2, random_state=42, stratify=y)
        kf = StratifiedKFold(n_splits=5, shuffle=True,
                             random_state=42)
        best_auc = 0
        for name, pipeline in self.ALGORITHMS.items():
            pipeline.fit(X_train, y_train)
            try:
                y_score = pipeline.decision_function(X_test)
            except AttributeError:
                y_score = pipeline.predict_proba(X_test)[:,1]
            auc = roc_auc_score(y_test, y_score)
            cv  = cross_val_score(
                pipeline, X, y, cv=kf, scoring="f1")
            logger.info(f"{name}: AUC={auc:.4f} "
                        f"CV-F1={cv.mean():.4f}+-{cv.std():.4f}")
            if auc > best_auc:
                best_auc = auc; self.best_model = pipeline
        joblib.dump(self.best_model, MODELS_DIR/"spam_clf.pkl")
\end{lstlisting}
\end{frame}

%%%%%%%%%%%%%%%%%%%%%%%%%%%%%%%%%%%%%%%%%%%%%%%%%%%%%%%%%%%%%%%%%%%%%%%%%%%%%%%%%%
\begin{frame}[fragile]\frametitle{Module 1: ML Code: HousePricePredictor}
\begin{lstlisting}
class HousePricePredictor:
    def __init__(self, data_path=DATA_DIR/"house_prices.csv"):
        self.data_path = data_path; self.best_model = None

    def feature_engineering(self, df):
        df = df.copy()
        df["total_rooms"]  = df["bedrooms"] + df["bathrooms"]
        df["amenity_score"]= df["has_garage"] + df["has_garden"]
        df["age_category"] = pd.cut(
            df["age_years"], bins=[0,5,15,30,100],
            labels=["new","recent","old","vintage"])
        df["age_category"] = LabelEncoder().fit_transform(
            df["age_category"])
        return df

    def run(self):
        df = self.feature_engineering(
            pd.read_csv(self.data_path))
        feats = ["area_sqft","bedrooms","bathrooms",
            "age_years","distance_center_km",
            "has_garage","has_garden","floor_level",
            "total_rooms","amenity_score","age_category"]
        X, y = df[feats], df["price_lakh"]
        Xtr,Xte,ytr,yte = train_test_split(
            X, y, test_size=0.2, random_state=42)
        for name, pipe in {
            "Ridge": Pipeline([("sc",StandardScaler()),
                               ("r",Ridge(alpha=1.0))]),
            "GBM":   Pipeline([("sc",StandardScaler()),
                ("r",GradientBoostingRegressor(
                    n_estimators=200,learning_rate=0.05))])
        }.items():
            pipe.fit(Xtr, ytr)
            logger.info(f"{name} R2={r2_score(yte,pipe.predict(Xte)):.4f}")
\end{lstlisting}
\end{frame}

%%%%%%%%%%%%%%%%%%%%%%%%%%%%%%%%%%%%%%%%%%%%%%%%%%%%%%%%%%%%%%%%%%%%%%%%%%%%%%%%%%
\begin{frame}[fragile]\frametitle{Module 1: ML Code: main() \& Expected Output}
\begin{lstlisting}
def main():
    # --- Task 1: Spam Classification ---
    clf = SpamClassifier()
    clf.run()
    test_msgs = [
        "WIN a FREE prize! Call 0800-123456 NOW!",
        "Hey, are you coming to lunch at 1pm?",
    ]
    for msg in test_msgs:
        pred = clf.best_model.predict([msg])[0]
        label = "SPAM" if pred == 1 else "HAM"
        logger.info(f"[{label}] {msg[:50]}")

    # --- Task 2: House Price Regression ---
    pred = HousePricePredictor()
    pred.run()
    sample = {"area_sqft":1200,"bedrooms":3,"bathrooms":2,
              "age_years":10,"distance_center_km":5.0,
              "has_garage":1,"has_garden":1,"floor_level":2}
    logger.info(f"Predicted price: {pred.best_model.predict(
        pd.DataFrame([sample]))[0]:.2f} Lakh")

if __name__ == "__main__":
    main()
\end{lstlisting}
\begin{block}{Expected terminal output (abridged)}
\texttt{Naive Bayes: AUC=0.9750 CV-F1=0.942$\pm$0.023}\\
\texttt{[SPAM] WIN a FREE prize! Call 0800-123456 NOW!}\\
\texttt{Predicted price: 72.45 Lakh}
\end{block}
\end{frame}

%%%%%%%%%%%%%%%%%%%%%%%%%%%%%%%%%%%%%%%%%%%%%%%%%%%%%%%%%%%%%%%%%%%%%%%%%%%%%%%%%%
\begin{frame}[fragile]\frametitle{}
\begin{center}
{\Large \textbf{Module 2: Deep Learning (PyTorch CPU)}}\

\texttt{src/dl/driver.py} $\cdot$ Bidirectional LSTM $\cdot$ Tabular MLP
\end{center}
\end{frame}

%%%%%%%%%%%%%%%%%%%%%%%%%%%%%%%%%%%%%%%%%%%%%%%%%%%%%%%%%%%%%%%%%%%%%%%%%%%%%%%%%%
\begin{frame}[fragile]\frametitle{Module 2: DL Concepts}
\begin{itemize}
  \item \textbf{Why PyTorch on CPU?} PyTorch ships a CPU-only build (\texttt{whl/cpu}) that is $\sim$2$\times$ smaller and installs cleanly without CUDA. All tensor operations and autograd work identically: just slower.
  \item \textbf{Bidirectional LSTM:} Each token attends to tokens both before and after it. The final hidden states of the forward and backward passes are concatenated, giving a context-aware sentence embedding. Better than unidirectional for classification.
  \item \textbf{Batch Normalisation:} Normalises layer inputs per mini-batch, reducing internal covariate shift. Allows higher learning rates and speeds convergence on tabular data.
  \item \textbf{Huber Loss:} Combines MSE (for small errors) and MAE (for large errors), making regression robust to outliers in the house price dataset.
  \item \textbf{CPU thread tuning:} \texttt{torch.set\_num\_threads(4)} pins PyTorch to physical cores; hyperthreads often slow matrix multiply on old CPUs.
\end{itemize}
\begin{block}{Run It}
  \texttt{cd src/dl \&\& python driver.py} $\quad$ then $\quad$ \texttt{make ui-dl}
\end{block}
\end{frame}

%%%%%%%%%%%%%%%%%%%%%%%%%%%%%%%%%%%%%%%%%%%%%%%%%%%%%%%%%%%%%%%%%%%%%%%%%%%%%%%%%%
\begin{frame}[fragile]\frametitle{Module 2: DL Code: Dataset \& LSTM Model}
\begin{lstlisting}
import torch, torch.nn as nn, torch.optim as optim
from torch.utils.data import Dataset, DataLoader, TensorDataset

torch.set_num_threads(4)   # physical cores only

class SentimentDataset(Dataset):
    TEXTS = [
        ("I love this product! Works perfectly.", 1),
        ("Terrible quality, broke after one day.", 0),
        # ... 22 more reviews (see driver.py) ...
    ]
    def __init__(self):
        self.vocab = {"<pad>":0,"<unk>":1}
        for text,_ in self.TEXTS:
            for w in text.lower().split():
                w = "".join(c for c in w if c.isalpha())
                if w and w not in self.vocab:
                    self.vocab[w] = len(self.vocab)
        self.max_len = 15
    def encode(self, text):
        ids = [self.vocab.get(
            "".join(c for c in t if c.isalpha()),1)
               for t in text.lower().split()]
        return ids[:self.max_len] + \
               [0]*max(0, self.max_len - len(ids))
    def __len__(self): return len(self.TEXTS)
    def __getitem__(self, i):
        t,l = self.TEXTS[i]
        return (torch.tensor(self.encode(t),dtype=torch.long),
                torch.tensor(l, dtype=torch.float))
\end{lstlisting}
\end{frame}

%%%%%%%%%%%%%%%%%%%%%%%%%%%%%%%%%%%%%%%%%%%%%%%%%%%%%%%%%%%%%%%%%%%%%%%%%%%%%%%%%%
\begin{frame}[fragile]\frametitle{Module 2: DL Code: LSTM Architecture}
\begin{lstlisting}
class LSTMClassifier(nn.Module):
    """Bidirectional LSTM for binary text classification."""
    def __init__(self, vocab_size, embed_dim=32,
                 hidden_dim=64, dropout=0.3):
        super().__init__()
        self.embedding = nn.Embedding(
            vocab_size, embed_dim, padding_idx=0)
        self.lstm = nn.LSTM(
            embed_dim, hidden_dim,
            batch_first=True, bidirectional=True,
            dropout=dropout)
        self.dropout = nn.Dropout(dropout)
        self.classifier = nn.Sequential(
            nn.Linear(hidden_dim * 2, 32),
            nn.ReLU(),
            nn.Dropout(dropout),
            nn.Linear(32, 1),
            nn.Sigmoid()
        )
    def forward(self, x):
        emb = self.dropout(self.embedding(x))
        _, (hidden, _) = self.lstm(emb)
        # concat forward + backward final hidden states
        ctx = torch.cat([hidden[-2], hidden[-1]], dim=1)
        return self.classifier(ctx).squeeze(1)
\end{lstlisting}
\end{frame}

%%%%%%%%%%%%%%%%%%%%%%%%%%%%%%%%%%%%%%%%%%%%%%%%%%%%%%%%%%%%%%%%%%%%%%%%%%%%%%%%%%
\begin{frame}[fragile]\frametitle{Module 2: DL Code: Training Loop}
\begin{lstlisting}
def train_lstm_classifier(epochs=30):
    dataset = SentimentDataset()
    loader  = DataLoader(dataset, batch_size=8, shuffle=True)
    model   = LSTMClassifier(vocab_size=len(dataset.vocab))
    opt     = optim.Adam(model.parameters(),
                         lr=1e-3, weight_decay=1e-4)
    loss_fn = nn.BCELoss()
    sched   = optim.lr_scheduler.StepLR(
                  opt, step_size=10, gamma=0.5)

    for epoch in range(epochs):
        model.train()
        epoch_loss, correct, total = 0.0, 0, 0
        for Xb, yb in loader:
            opt.zero_grad()
            preds = model(Xb)
            loss  = loss_fn(preds, yb)
            loss.backward()
            nn.utils.clip_grad_norm_(model.parameters(), 1.0)
            opt.step()
            epoch_loss += loss.item()
            correct += ((preds>0.5)==yb.bool()).sum().item()
            total   += len(yb)
        sched.step()
        if (epoch+1) % 10 == 0:
            logger.info(f"Epoch {epoch+1} "
                f"Loss={epoch_loss/len(loader):.4f} "
                f"Acc={correct/total:.3f}")
    torch.save(model.state_dict(), "models/lstm.pt")
    return model
\end{lstlisting}
\end{frame}

%%%%%%%%%%%%%%%%%%%%%%%%%%%%%%%%%%%%%%%%%%%%%%%%%%%%%%%%%%%%%%%%%%%%%%%%%%%%%%%%%%
\begin{frame}[fragile]\frametitle{Module 2: DL Code: Tabular MLP}
\begin{lstlisting}
class TabularMLP(nn.Module):
    """MLP with BatchNorm for house price regression."""
    def __init__(self, input_dim):
        super().__init__()
        self.network = nn.Sequential(
            nn.Linear(input_dim, 128),
            nn.BatchNorm1d(128), nn.ReLU(), nn.Dropout(0.2),
            nn.Linear(128, 64),
            nn.BatchNorm1d(64),  nn.ReLU(), nn.Dropout(0.2),
            nn.Linear(64, 32),   nn.ReLU(),
            nn.Linear(32, 1)
        )
    def forward(self, x):
        return self.network(x).squeeze(1)

def train_tabular_mlp(epochs=100):
    # load house_prices.csv, StandardScaler, split
    # ... (see full driver.py)
    model = TabularMLP(input_dim=8)
    opt   = optim.AdamW(model.parameters(),
                        lr=1e-3, weight_decay=1e-3)
    for epoch in range(epochs):
        model.train()
        for Xb, yb in loader:
            opt.zero_grad()
            loss = nn.HuberLoss()(model(Xb), yb)
            loss.backward(); opt.step()
    # compute R-squared on test set
    # logger.info(f"Test R2: {r2:.4f}")
    torch.save(model.state_dict(), "models/mlp.pt")
\end{lstlisting}
\end{frame}

%%%%%%%%%%%%%%%%%%%%%%%%%%%%%%%%%%%%%%%%%%%%%%%%%%%%%%%%%%%%%%%%%%%%%%%%%%%%%%%%%%
\begin{frame}[fragile]\frametitle{}
\begin{center}
{\Large \textbf{Module 3: RAG Pipeline}}\

\texttt{src/rag/driver.py} $\cdot$ LangChain $\cdot$ ChromaDB $\cdot$ Ollama
\end{center}
\end{frame}

%%%%%%%%%%%%%%%%%%%%%%%%%%%%%%%%%%%%%%%%%%%%%%%%%%%%%%%%%%%%%%%%%%%%%%%%%%%%%%%%%%
\begin{frame}[fragile]\frametitle{Module 3: RAG Concepts}
\begin{itemize}
  \item \textbf{RAG = Retrieve-then-Generate:} Instead of asking an LLM to answer from memory, we first retrieve the most relevant passages from a document store and include them in the prompt as context.
  \item \textbf{Embeddings:} \texttt{nomic-embed-text} converts text chunks to 768-dimensional vectors (via Ollama, CPU). Semantically similar texts have nearby vectors.
  \item \textbf{ChromaDB:} A lightweight vector database that stores embeddings on disk (\texttt{.chroma\_db/}). At query time, it computes cosine similarity to find the top-K most relevant chunks in milliseconds.
  \item \textbf{Chunking:} \texttt{RecursiveCharacterTextSplitter} splits documents into 500-character chunks with 50-character overlap so that no important sentence is cut at a boundary.
  \item \textbf{RetrievalQA Chain:} LangChain's \texttt{RetrievalQA} wires retriever $\rightarrow$ prompt template $\rightarrow$ Ollama LLM into one callable.
  \item \textbf{Offline:} Index is built once and persisted. Subsequent runs skip re-embedding.
\end{itemize}
\begin{block}{Run It}
  Start Ollama first: \texttt{ollama serve \&} $\quad$ then \texttt{make run-rag}
\end{block}
\end{frame}

%%%%%%%%%%%%%%%%%%%%%%%%%%%%%%%%%%%%%%%%%%%%%%%%%%%%%%%%%%%%%%%%%%%%%%%%%%%%%%%%%%
\begin{frame}[fragile]\frametitle{Module 3: RAG Code: Config \& Prompt}
\begin{lstlisting}
from langchain.text_splitter import RecursiveCharacterTextSplitter
from langchain_community.document_loaders import TextLoader
from langchain_community.vectorstores import Chroma
from langchain_community.embeddings import OllamaEmbeddings
from langchain_community.llms import Ollama
from langchain.chains import RetrievalQA
from langchain.prompts import PromptTemplate
from langchain.schema import Document

EMBED_MODEL  = "nomic-embed-text"
LLM_MODEL    = "gemma2:2b"
CHUNK_SIZE   = 500
CHUNK_OVERLAP= 50
TOP_K        = 3

RAG_PROMPT = PromptTemplate(
    input_variables=["context", "question"],
    template="""Answer ONLY from the context below.
If unsure, say "I don't have that information."

Context:
{context}

Question: {question}

Answer:""")
\end{lstlisting}
\end{frame}

%%%%%%%%%%%%%%%%%%%%%%%%%%%%%%%%%%%%%%%%%%%%%%%%%%%%%%%%%%%%%%%%%%%%%%%%%%%%%%%%%%
\begin{frame}[fragile]\frametitle{Module 3: RAG Code: Build Index}
\begin{lstlisting}
class LocalRAGPipeline:
    def __init__(self, embed_model=EMBED_MODEL,
                 llm_model=LLM_MODEL,
                 kb_dir=Path("knowledge_base"),
                 persist_dir=Path(".chroma_db")):
        self.embed_model  = embed_model
        self.llm_model    = llm_model
        self.kb_dir       = kb_dir
        self.persist_dir  = persist_dir
        self.vectorstore  = None
        self.qa_chain     = None

    def build_index(self, force_rebuild=False):
        db_file = self.persist_dir / "chroma.sqlite3"
        embeddings = OllamaEmbeddings(model=self.embed_model)
        if not force_rebuild and db_file.exists():
            self.vectorstore = Chroma(
                persist_directory=str(self.persist_dir),
                embedding_function=embeddings)
            return
        docs = [TextLoader(str(f), encoding="utf-8").load()[0]
                for f in self.kb_dir.glob("*.txt")]
        splitter = RecursiveCharacterTextSplitter(
            chunk_size=CHUNK_SIZE, chunk_overlap=CHUNK_OVERLAP)
        chunks = splitter.split_documents(docs)
        self.vectorstore = Chroma.from_documents(
            chunks, embeddings,
            persist_directory=str(self.persist_dir))
        self.vectorstore.persist()
        logger.info(f"Indexed {len(chunks)} chunks")
\end{lstlisting}
\end{frame}

%%%%%%%%%%%%%%%%%%%%%%%%%%%%%%%%%%%%%%%%%%%%%%%%%%%%%%%%%%%%%%%%%%%%%%%%%%%%%%%%%%
\begin{frame}[fragile]\frametitle{Module 3: RAG Code: Build Chain \& Query}
\begin{lstlisting}
    def build_chain(self):
        llm = Ollama(model=self.llm_model,
                     temperature=0.1, num_predict=512)
        retriever = self.vectorstore.as_retriever(
            search_type="similarity",
            search_kwargs={"k": TOP_K})
        self.qa_chain = RetrievalQA.from_chain_type(
            llm=llm,
            chain_type="stuff",
            retriever=retriever,
            chain_type_kwargs={"prompt": RAG_PROMPT},
            return_source_documents=True)

    def query(self, question: str) -> dict:
        import time
        t0 = time.time()
        result = self.qa_chain({"query": question})
        return {
            "answer":  result["result"],
            "sources": [d.metadata.get("source","?")
                        for d in result["source_documents"]],
            "latency": round(time.time()-t0, 2)
        }

# Usage
pipeline = LocalRAGPipeline()
pipeline.build_index()
pipeline.build_chain()
r = pipeline.query("Which model is best for multilingual use?")
print(r["answer"], r["latency"])
\end{lstlisting}
\end{frame}

%%%%%%%%%%%%%%%%%%%%%%%%%%%%%%%%%%%%%%%%%%%%%%%%%%%%%%%%%%%%%%%%%%%%%%%%%%%%%%%%%%
\begin{frame}[fragile]\frametitle{}
\begin{center}
{\Large \textbf{Module 4: AI Agents with LangGraph}}\

\texttt{src/agents/driver.py} $\cdot$ State machines $\cdot$ Local LLM tools
\end{center}
\end{frame}

%%%%%%%%%%%%%%%%%%%%%%%%%%%%%%%%%%%%%%%%%%%%%%%%%%%%%%%%%%%%%%%%%%%%%%%%%%%%%%%%%%
\begin{frame}[fragile]\frametitle{Module 4: Agent Concepts}
\begin{itemize}
  \item \textbf{LangGraph:} Builds agents as typed state machines (directed graphs). Each node is a Python function that reads/writes a shared \texttt{TypedDict} state. Edges define the flow: fixed or conditional.
  \item \textbf{ResearchAgent: Linear Graph:} \texttt{START} $\to$ \texttt{plan} $\to$ \texttt{research} $\to$ \texttt{synthesize} $\to$ \texttt{END}. Each node calls the local Ollama LLM with a focused prompt.
  \item \textbf{CodeReviewAgent: Conditional Graph:} After \texttt{syntax\_check}, a router reads the result and branches to \texttt{pass\_review}, \texttt{minor\_review}, or \texttt{major\_review}.
  \item \textbf{Tools:} LangChain \texttt{@tool} decorators expose Python functions (calculator, word counter, AST syntax checker) that agents can call without an LLM round-trip.
  \item \textbf{No OpenAI API key needed:} All LLM calls go to \texttt{localhost:11434} (Ollama).
\end{itemize}
\begin{block}{Run It}
  \texttt{make run-agents} $\quad$ then $\quad$ \texttt{make ui-agents}
\end{block}
\end{frame}

%%%%%%%%%%%%%%%%%%%%%%%%%%%%%%%%%%%%%%%%%%%%%%%%%%%%%%%%%%%%%%%%%%%%%%%%%%%%%%%%%%
\begin{frame}[fragile]\frametitle{Module 4: Agent Code: Tools}
\begin{lstlisting}
import math, json
from langchain.tools import tool
from langgraph.graph import StateGraph, END
from langchain_community.llms import Ollama
from typing import TypedDict

@tool
def calculator(expression: str) -> str:
    """Safely evaluate a math expression."""
    allowed = {
        "abs":abs,"round":round,"sqrt":math.sqrt,
        "pow":pow,"log":math.log,"pi":math.pi,"e":math.e,
        "sin":math.sin,"cos":math.cos,"tan":math.tan,
    }
    try:
        result = eval(expression,{"__builtins__":{}},allowed)
        return f"Result: {result}"
    except Exception as ex:
        return f"Error: {ex}"

@tool
def python_syntax_checker(code: str) -> str:
    """Check Python code for syntax errors via ast.parse."""
    import ast
    try:
        tree = ast.parse(code)
        return json.dumps({
            "status":"valid",
            "functions": sum(1 for n in ast.walk(tree)
                             if isinstance(n,ast.FunctionDef)),
            "lines": len(code.splitlines())})
    except SyntaxError as e:
        return json.dumps({"status":"error","msg":str(e),
                           "line":e.lineno})
\end{lstlisting}
\end{frame}

%%%%%%%%%%%%%%%%%%%%%%%%%%%%%%%%%%%%%%%%%%%%%%%%%%%%%%%%%%%%%%%%%%%%%%%%%%%%%%%%%%
\begin{frame}[fragile]\frametitle{Module 4: Agent Code: ResearchAgent Graph}
\begin{lstlisting}
class ResearchState(TypedDict):
    topic: str;  research_plan: str
    gathered_info: str; final_report: str; step_count: int

class ResearchAgent:
    def __init__(self, model="gemma2:2b"):
        self.llm = Ollama(model=model,
                          temperature=0.3,num_predict=400)
        self.graph = self._build_graph()

    def _plan(self, s: ResearchState) -> ResearchState:
        plan = self.llm.invoke(
            f"Write a 3-point research plan for: {s['topic']}")
        return {**s,"research_plan":plan,"step_count":s["step_count"]+1}

    def _research(self, s: ResearchState) -> ResearchState:
        info = self.llm.invoke(
            f"Topic: {s['topic']}\nPlan: {s['research_plan']}\n"
            "Provide 3 key facts. Under 150 words.")
        return {**s,"gathered_info":info,"step_count":s["step_count"]+1}

    def _synthesize(self, s: ResearchState) -> ResearchState:
        rep = self.llm.invoke(
            f"Synthesize into a 2-paragraph report:\n{s['gathered_info']}")
        return {**s,"final_report":rep,"step_count":s["step_count"]+1}

    def _build_graph(self):
        wf = StateGraph(ResearchState)
        wf.add_node("plan",       self._plan)
        wf.add_node("research",   self._research)
        wf.add_node("synthesize", self._synthesize)
        wf.set_entry_point("plan")
        wf.add_edge("plan","research")
        wf.add_edge("research","synthesize")
        wf.add_edge("synthesize", END)
        return wf.compile()
\end{lstlisting}
\end{frame}

%%%%%%%%%%%%%%%%%%%%%%%%%%%%%%%%%%%%%%%%%%%%%%%%%%%%%%%%%%%%%%%%%%%%%%%%%%%%%%%%%%
\begin{frame}[fragile]\frametitle{Module 4: Agent Code: CodeReview Conditional Graph}
\begin{lstlisting}
class CodeReviewState(TypedDict):
    code: str; syntax_check: str
    severity: str; review_feedback: str; step_count: int

class CodeReviewAgent:
    def __init__(self, model="gemma2:2b"):
        self.llm = Ollama(model=model,temperature=0.2)
        self.graph = self._build_graph()

    def _syntax(self, s):
        return {**s,"syntax_check":python_syntax_checker.invoke(s["code"]),
                "step_count":s["step_count"]+1}

    def _route(self, s) -> str:
        data = json.loads(s["syntax_check"])
        if data["status"]=="syntax_error": return "major"
        if data.get("functions",0)==0:     return "minor"
        return "pass"

    def _pass_review(self, s):
        fb = self.llm.invoke(f"Good code, 2-sentence review:\n{s['code'][:400]}")
        return {**s,"severity":"pass","review_feedback":fb}

    def _major_review(self, s):
        fb = self.llm.invoke(f"Syntax error {s['syntax_check']}.\nFix:\n{s['code'][:400]}")
        return {**s,"severity":"major","review_feedback":fb}

    def _build_graph(self):
        wf = StateGraph(CodeReviewState)
        for name,fn in [("chk",self._syntax),
                        ("pass",self._pass_review),
                        ("major",self._major_review)]:
            wf.add_node(name, fn)
        wf.set_entry_point("chk")
        wf.add_conditional_edges("chk", self._route,
            {"pass":"pass","minor":"pass","major":"major"})
        for n in ["pass","major"]: wf.add_edge(n, END)
        return wf.compile()
\end{lstlisting}
\end{frame}

%%%%%%%%%%%%%%%%%%%%%%%%%%%%%%%%%%%%%%%%%%%%%%%%%%%%%%%%%%%%%%%%%%%%%%%%%%%%%%%%%%
\begin{frame}[fragile]\frametitle{}
\begin{center}
{\Large \textbf{Module 5: Computer Vision (OpenCV)}}\

\texttt{src/cv/driver.py} $\cdot$ Haar Cascades $\cdot$ MOG2 Background Subtraction
\end{center}
\end{frame}

%%%%%%%%%%%%%%%%%%%%%%%%%%%%%%%%%%%%%%%%%%%%%%%%%%%%%%%%%%%%%%%%%%%%%%%%%%%%%%%%%%
\begin{frame}[fragile]\frametitle{Module 5: CV Concepts}
\begin{itemize}
  \item \textbf{Haar Cascade:} A sliding-window classifier trained on thousands of face/non-face images. Uses Haar-like features (differences of rectangular sums). Cascade architecture rejects non-face regions early: very fast on CPU ($\sim$15\,FPS).
  \item \textbf{MOG2 (Mixture of Gaussians v2):} Models each pixel's intensity as a mixture of Gaussians learned from recent frames. Pixels deviating from the background distribution are flagged as foreground (motion). Shadow detection built-in.
  \item \textbf{Canny Edge Detection:} Two-threshold hysteresis edge detector. Pixels above \texttt{threshold2} are strong edges; pixels between the thresholds are kept only if connected to a strong edge.
  \item \textbf{Laplacian Variance (sharpness):} The variance of the Laplacian of an image measures the spread of edge strengths. Low variance $\Rightarrow$ blurry image.
  \item \textbf{No GPU, no DL:} All algorithms are classical signal processing: run at full speed on a 2012 i3 CPU.
\end{itemize}
\begin{block}{Run It}
  \texttt{make run-cv} $\quad$ then $\quad$ \texttt{streamlit run src/cv/ui/app.py}
\end{block}
\end{frame}

%%%%%%%%%%%%%%%%%%%%%%%%%%%%%%%%%%%%%%%%%%%%%%%%%%%%%%%%%%%%%%%%%%%%%%%%%%%%%%%%%%
\begin{frame}[fragile]\frametitle{Module 5: CV Code: Face Detector}
\begin{lstlisting}
import cv2
import numpy as np
from pathlib import Path

MODELS_DIR = Path(__file__).parent / "models"

class HaarFaceDetector:
    def __init__(self):
        cascade = MODELS_DIR/"haarcascade_frontalface_default.xml"
        if not cascade.exists():
            # fall back to OpenCV bundled cascades
            cascade = Path(cv2.__file__).parent/"data"/cascade.name
        self.face_cascade = cv2.CascadeClassifier(str(cascade))

    def detect(self, image: np.ndarray):
        gray = cv2.cvtColor(image, cv2.COLOR_BGR2GRAY) \
               if image.ndim == 3 else image
        gray = cv2.equalizeHist(gray)
        faces = self.face_cascade.detectMultiScale(
            gray, scaleFactor=1.1, minNeighbors=5,
            minSize=(30, 30), flags=cv2.CASCADE_SCALE_IMAGE)
        return list(faces) if len(faces) > 0 else []

    def annotate(self, image, faces):
        out = image.copy()
        for (x, y, w, h) in faces:
            cv2.rectangle(out,(x,y),(x+w,y+h),(0,255,0),2)
            cv2.putText(out,"Face",(x,y-8),
                cv2.FONT_HERSHEY_SIMPLEX,0.6,(0,255,0),2)
        return out
\end{lstlisting}
\end{frame}

%%%%%%%%%%%%%%%%%%%%%%%%%%%%%%%%%%%%%%%%%%%%%%%%%%%%%%%%%%%%%%%%%%%%%%%%%%%%%%%%%%
\begin{frame}[fragile]\frametitle{Module 5: CV Code: Motion Detector \& Image Analyser}
\begin{lstlisting}
class MotionDetector:
    def __init__(self, history=500, threshold=16.0):
        self.bg = cv2.createBackgroundSubtractorMOG2(
            history=history, varThreshold=threshold,
            detectShadows=True)
        self.kernel = cv2.getStructuringElement(
            cv2.MORPH_ELLIPSE, (5, 5))

    def process_frame(self, frame):
        mask = self.bg.apply(frame)
        _, mask = cv2.threshold(mask, 200, 255,
                                cv2.THRESH_BINARY)
        mask = cv2.morphologyEx(mask,cv2.MORPH_OPEN,self.kernel)
        mask = cv2.dilate(mask, self.kernel, iterations=2)
        contours,_ = cv2.findContours(
            mask, cv2.RETR_EXTERNAL, cv2.CHAIN_APPROX_SIMPLE)
        moving = [c for c in contours
                  if cv2.contourArea(c) > 500]
        ratio  = sum(cv2.contourArea(c) for c in moving) \
                 / (frame.shape[0]*frame.shape[1])
        return mask, moving, ratio

class ImageAnalyzer:
    def analyze(self, image):
        gray = cv2.cvtColor(image,cv2.COLOR_BGR2GRAY) \
               if image.ndim==3 else image
        edges   = cv2.Canny(gray, 50, 150)
        lap_var = cv2.Laplacian(gray, cv2.CV_64F).var()
        hist    = cv2.calcHist([gray],[0],None,[256],[0,256])
        return {"sharpness": round(float(lap_var),2),
                "edge_density":
                    round(float(edges.sum()/(255*gray.size)),4),
                "is_blurry": lap_var < 100}
\end{lstlisting}
\end{frame}

%%%%%%%%%%%%%%%%%%%%%%%%%%%%%%%%%%%%%%%%%%%%%%%%%%%%%%%%%%%%%%%%%%%%%%%%%%%%%%%%%%
\begin{frame}[fragile]\frametitle{}
\begin{center}
{\Large \textbf{Module 6: Local LLM Inference \& Benchmarking}}\

\texttt{src/llm\_inference/driver.py} $\cdot$ Ollama REST API $\cdot$ Prompt Engineering
\end{center}
\end{frame}

%%%%%%%%%%%%%%%%%%%%%%%%%%%%%%%%%%%%%%%%%%%%%%%%%%%%%%%%%%%%%%%%%%%%%%%%%%%%%%%%%%
\begin{frame}[fragile]\frametitle{Module 6: LLM Inference Concepts}
\begin{itemize}
  \item \textbf{Ollama REST API:} Ollama exposes \texttt{http://localhost:11434} with endpoints: \texttt{/api/generate} (completion), \texttt{/api/chat} (multi-turn), \texttt{/api/embeddings}, \texttt{/api/tags} (list models).
  \item \textbf{Tokens per second (tok/s):} The primary CPU inference speed metric. On i3-8130U: Gemma2:2b $\approx$ 8\,tok/s, Qwen2:1.5b $\approx$ 12\,tok/s, TinyLlama $\approx$ 15\,tok/s.
  \item \textbf{Prompt Engineering Patterns:}
  \begin{itemize}
    \item \textit{Zero-shot:} Direct question, no examples
    \item \textit{Few-shot:} Provide 2--3 input/output examples before the query
    \item \textit{Chain-of-thought:} Ask the model to reason step-by-step before answering
    \item \textit{Structured output:} Instruct JSON-only response, parse with \texttt{json.loads}
  \end{itemize}
  \item \textbf{Temperature:} 0.0--0.2 for deterministic tasks (code, JSON); 0.7--0.9 for creative tasks.
\end{itemize}
\begin{block}{Run It}
  \texttt{make run-llm} $\quad$ then $\quad$ \texttt{make ui-llm}
\end{block}
\end{frame}

%%%%%%%%%%%%%%%%%%%%%%%%%%%%%%%%%%%%%%%%%%%%%%%%%%%%%%%%%%%%%%%%%%%%%%%%%%%%%%%%%%
\begin{frame}[fragile]\frametitle{Module 6: LLM Code: Ollama Client}
\begin{lstlisting}
import httpx, time, json

OLLAMA_BASE = "http://localhost:11434"

class OllamaClient:
    def __init__(self):
        self.client = httpx.Client(timeout=120.0)

    def is_running(self) -> bool:
        try:
            return self.client.get(
                f"{OLLAMA_BASE}/api/tags").status_code==200
        except Exception: return False

    def list_models(self) -> list:
        r = self.client.get(f"{OLLAMA_BASE}/api/tags")
        return r.json().get("models", [])

    def generate(self, model, prompt,
                 temperature=0.7, max_tokens=512) -> dict:
        payload = {"model":model, "prompt":prompt,
                   "stream":False,
                   "options":{"temperature":temperature,
                               "num_predict":max_tokens}}
        t0 = time.time()
        r  = self.client.post(
                 f"{OLLAMA_BASE}/api/generate", json=payload)
        elapsed = time.time()-t0
        data = r.json()
        return {"response": data.get("response",""),
                "latency_s": round(elapsed,2),
                "tok_per_s": round(
                    data.get("eval_count",0)/max(elapsed,0.001),1)}
\end{lstlisting}
\end{frame}

%%%%%%%%%%%%%%%%%%%%%%%%%%%%%%%%%%%%%%%%%%%%%%%%%%%%%%%%%%%%%%%%%%%%%%%%%%%%%%%%%%
\begin{frame}[fragile]\frametitle{Module 6: LLM Code: Benchmarker \& Prompt Patterns}
\begin{lstlisting}
PROMPTS = [
    "Explain attention mechanism in 2 sentences.",
    "Write a Python function to check if n is prime.",
    "List 5 Linux commands for system monitoring.",
]

class ModelBenchmarker:
    def __init__(self):
        self.client = OllamaClient()
        self.models = [m["name"]
                       for m in self.client.list_models()]

    def run(self, n_prompts=3):
        for model in self.models:
            times = []
            for prompt in PROMPTS[:n_prompts]:
                r = self.client.generate(model, prompt,
                                         max_tokens=200)
                times.append(r["tok_per_s"])
                logger.info(f"{model}: {r['tok_per_s']} tok/s")
            logger.info(f"AVG {model}: "
                        f"{sum(times)/len(times):.1f} tok/s")

class PromptEngineer:
    def __init__(self, model="gemma2:2b"):
        self.c = OllamaClient(); self.model = model

    def chain_of_thought(self, problem: str) -> str:
        prompt = (f"Solve step by step:\n{problem}\n\n"
                  "Step 1:")
        return self.c.generate(self.model, prompt,
                                temperature=0.3)["response"]

    def structured_json(self, task, schema) -> dict:
        prompt = f"{task}\nRespond ONLY as JSON:\n{schema}"
        r = self.c.generate(self.model, prompt, temperature=0.1)
        return json.loads(r["response"].strip())
\end{lstlisting}
\end{frame}

%%%%%%%%%%%%%%%%%%%%%%%%%%%%%%%%%%%%%%%%%%%%%%%%%%%%%%%%%%%%%%%%%%%%%%%%%%%%%%%%%%
\begin{frame}[fragile]\frametitle{}
\begin{center}
{\Large \textbf{Integration Tests \& Running the Full Suite}}\

\texttt{tests/test\_integration.py} $\cdot$ pytest $\cdot$ No Ollama required
\end{center}
\end{frame}

%%%%%%%%%%%%%%%%%%%%%%%%%%%%%%%%%%%%%%%%%%%%%%%%%%%%%%%%%%%%%%%%%%%%%%%%%%%%%%%%%%
\begin{frame}[fragile]\frametitle{Tests: What Is Checked}
\begin{itemize}
  \item \textbf{TestRepoStructure:} Parametrised: asserts every directory and file listed in the README actually exists on disk. Catches copy-paste errors immediately.
  \item \textbf{TestDatasets:} Loads both CSVs with pandas, checks row counts, column names, no null values, both class labels present.
  \item \textbf{TestImports:} Verifies \texttt{numpy, pandas, sklearn, torch, cv2, streamlit} are importable. Optional packages (\texttt{langchain, chromadb, ollama}) are skipped gracefully if absent.
  \item \textbf{TestCrossModule:} Imports each \texttt{driver.py} and verifies key classes/functions are callable without running training.
  \item \textbf{TestOfflineReadiness:} Scans driver files for hardcoded API key patterns (\texttt{sk-...}), checks \texttt{requirements.txt} has $\geq$10 packages, verifies Makefile targets exist.
\end{itemize}
\begin{block}{Module-Level Tests}
Each module also has its own \texttt{tests/test\_*.py}: \texttt{TestSpamClassifier}, \texttt{TestHousePricePredictor}, \texttt{TestLSTMClassifier}, \texttt{TestCalculatorTool}, etc.
\end{block}
\end{frame}

%%%%%%%%%%%%%%%%%%%%%%%%%%%%%%%%%%%%%%%%%%%%%%%%%%%%%%%%%%%%%%%%%%%%%%%%%%%%%%%%%%
\begin{frame}[fragile]\frametitle{Tests: Running \& Reading Output}
\begin{lstlisting}
conda activate aiooll

# Run all integration tests (no Ollama needed)
pytest tests/ -v --timeout=60

# Run module-level ML tests
pytest src/ml/tests/ -v --timeout=120

# Run DL tests (slow tests need flag)
pytest src/dl/tests/ -v -m "not slow" --timeout=60

# Run agent tool tests (no LLM needed)
pytest src/agents/tests/ -v --timeout=30

# Run everything at once
make test

# Check coverage (optional)
pytest tests/ src/ml/tests/ --cov=src \
       --cov-report=term-missing
\end{lstlisting}
\begin{block}{Expected output (abridged)}
\texttt{PASSED tests/test\_integration.py::TestRepoStructure::...}\\
\texttt{PASSED tests/test\_integration.py::TestDatasets::test\_sms\_spam\_csv\_loads}\\
\texttt{PASSED src/ml/tests/test\_ml.py::TestSpamClassifier::test\_auc\_threshold}
\end{block}
\end{frame}

%%%%%%%%%%%%%%%%%%%%%%%%%%%%%%%%%%%%%%%%%%%%%%%%%%%%%%%%%%%%%%%%%%%%%%%%%%%%%%%%%%
\begin{frame}[fragile]\frametitle{Tests: Integration Test Code (Part 1)}
\begin{lstlisting}
import pytest, importlib
from pathlib import Path
ROOT = Path(__file__).parent.parent

REQUIRED_FILES = [
    "README.md","requirements.txt","Makefile",
    "src/ml/driver.py","src/ml/ui/app.py",
    "src/ml/tests/test_ml.py",
    "src/ml/data/sms_spam.csv",
    "src/dl/driver.py","src/rag/driver.py",
    "src/agents/driver.py","src/cv/driver.py",
    "src/llm_inference/driver.py",
]

class TestRepoStructure:
    @pytest.mark.parametrize("fp", REQUIRED_FILES)
    def test_file_exists(self, fp):
        assert (ROOT/fp).is_file(), f"Missing: {fp}"

class TestDatasets:
    def test_sms_spam_loads(self):
        import pandas as pd
        df = pd.read_csv(ROOT/"src/ml/data/sms_spam.csv")
        assert len(df) >= 50
        assert set(df["label"].unique()) <= {"spam","ham"}

    def test_house_prices_no_nulls(self):
        import pandas as pd
        df = pd.read_csv(ROOT/"src/ml/data/house_prices.csv")
        assert df.isnull().sum().sum() == 0
        assert (df["price_lakh"] > 0).all()
\end{lstlisting}
\end{frame}

%%%%%%%%%%%%%%%%%%%%%%%%%%%%%%%%%%%%%%%%%%%%%%%%%%%%%%%%%%%%%%%%%%%%%%%%%%%%%%%%%%
\begin{frame}[fragile]\frametitle{Tests: Integration Test Code (Part 2)}
\begin{lstlisting}
CORE_PKGS = ["numpy","pandas","sklearn","torch","cv2",
             "streamlit","loguru"]

class TestImports:
    @pytest.mark.parametrize("pkg", CORE_PKGS)
    def test_importable(self, pkg):
        importlib.import_module(pkg)

    def test_torch_cpu(self):
        import torch
        x = torch.tensor([1.0, 2.0, 3.0])
        assert x.device.type == "cpu"

class TestOfflineReadiness:
    def test_no_hardcoded_api_keys(self):
        import re
        pattern = re.compile(r"sk-[A-Za-z0-9]{20,}")
        for driver in ROOT.glob("src/*/driver.py"):
            assert not pattern.search(driver.read_text()), \
                f"API key in {driver}"

    def test_makefile_targets(self):
        mf = (ROOT/"Makefile").read_text()
        for t in ["setup","test","run-ml","run-rag","clean"]:
            assert f"{t}:" in mf
\end{lstlisting}
\end{frame}

%%%%%%%%%%%%%%%%%%%%%%%%%%%%%%%%%%%%%%%%%%%%%%%%%%%%%%%%%%%%%%%%%%%%%%%%%%%%%%%%%%
\begin{frame}[fragile]\frametitle{}
\begin{center}
{\Large \textbf{Streamlit UIs: Launching \& Using}}\

\textit{Every module has an interactive dashboard}
\end{center}
\end{frame}

%%%%%%%%%%%%%%%%%%%%%%%%%%%%%%%%%%%%%%%%%%%%%%%%%%%%%%%%%%%%%%%%%%%%%%%%%%%%%%%%%%
\begin{frame}[fragile]\frametitle{Streamlit UI: Launch All Dashboards}
\begin{lstlisting}
conda activate aiooll

# ML Dashboard: spam detector + house price predictor
streamlit run src/ml/ui/app.py
# Opens: http://localhost:8501

# Deep Learning: training visualizer
streamlit run src/dl/ui/app.py
# Opens: http://localhost:8501

# RAG: document Q&A chat interface
# (requires: ollama serve & running)
streamlit run src/rag/ui/app.py

# AI Agents: thought-process visualizer
streamlit run src/agents/ui/app.py

# Computer Vision: face & motion detection
streamlit run src/cv/ui/app.py

# LLM Inference: benchmark + chat + prompt patterns
streamlit run src/llm_inference/ui/app.py
\end{lstlisting}
\begin{alertblock}{One UI at a time}
Each \texttt{streamlit run} command occupies port 8501. Run them in separate terminals, or stop one (\textbf{Ctrl+C}) before starting another.
\end{alertblock}
\end{frame}

%%%%%%%%%%%%%%%%%%%%%%%%%%%%%%%%%%%%%%%%%%%%%%%%%%%%%%%%%%%%%%%%%%%%%%%%%%%%%%%%%%
\begin{frame}[fragile]\frametitle{}
\begin{center}
{\Large \textbf{Hardware Optimisation for AI Workloads}}\

\textit{Squeeze every tok/s from your old i3}
\end{center}
\end{frame}

%%%%%%%%%%%%%%%%%%%%%%%%%%%%%%%%%%%%%%%%%%%%%%%%%%%%%%%%%%%%%%%%%%%%%%%%%%%%%%%%%%
\begin{frame}[fragile]\frametitle{Hardware Optimisation: Quick Wins}
\begin{lstlisting}
# 1. Set CPU governor to performance mode
echo performance | sudo tee \
    /sys/devices/system/cpu/cpu*/cpufreq/scaling_governor
cat /sys/devices/system/cpu/cpu0/cpufreq/scaling_governor
# should print: performance

# 2. Monitor CPU temperature (throttling check)
watch -n 2 sensors   # look for Package temp < 85 C

# 3. Check available RAM before inference
free -h              # aim for > 2 GB free

# 4. Drop file caches to free RAM
sudo sh -c "echo 3 > /proc/sys/vm/drop_caches"

# 5. Pin Ollama to physical cores only (2 on i3)
OMP_NUM_THREADS=2 ollama serve

# 6. Live monitoring during inference
btop                 # CPU, RAM, per-core load
\end{lstlisting}
\begin{block}{Model Selection by Available RAM}
3\,GB free: TinyLlama 1.1B $\cdot$ 4\,GB: Qwen2 1.5B $\cdot$ 6\,GB: Gemma2 2B $\cdot$ 8\,GB: Phi-3 Mini 3.8B
\end{block}
\end{frame}

%%%%%%%%%%%%%%%%%%%%%%%%%%%%%%%%%%%%%%%%%%%%%%%%%%%%%%%%%%%%%%%%%%%%%%%%%%%%%%%%%%
\begin{frame}[fragile]\frametitle{Troubleshooting Common Issues}
\begin{itemize}
  \item \textbf{``Ollama not running'':}
    \texttt{ollama serve \&} $\quad$ then wait 3 seconds and retry.
  \item \textbf{``Model not found'':}
    \texttt{ollama pull gemma2:2b} $\quad$ check disk space with \texttt{df -h}.
  \item \textbf{ParserError on sms\_spam.csv:}
    Use the fixed CSV from the repo (commas in text were causing 3-field parse). Driver now uses \texttt{quotechar='"'}.
  \item \textbf{PyTorch CUDA error:}
    Re-install CPU build: \texttt{pip install torch --index-url https://download.pytorch.org/whl/cpu}
  \item \textbf{Streamlit ``module not found'':}
    Ensure conda env is active: \texttt{conda activate aiooll}.
  \item \textbf{RAM OOM during inference:}
    Close browser tabs, run \texttt{sudo sh -c "echo 3 > /proc/sys/vm/drop\_caches"}, switch to a smaller model.
  \item \textbf{CV cascade not found:}
    Run \texttt{bash scripts/setup.sh}: it downloads Haar cascade XMLs automatically.
  \item \textbf{conda: command not found:}
    Run \texttt{source \$HOME/miniconda3/bin/activate} then \texttt{conda init bash \&\& source \textasciitilde/.bashrc}.
\end{itemize}
\end{frame}

%%%%%%%%%%%%%%%%%%%%%%%%%%%%%%%%%%%%%%%%%%%%%%%%%%%%%%%%%%%%%%%%%%%%%%%%%%%%%%%%%%
\begin{frame}[fragile]\frametitle{Workshop Summary \& Next Steps}
\begin{columns}
\column{0.5\textwidth}
\textbf{What you built today:}
\begin{itemize}
  \item[$\checkmark$] Xubuntu on old hardware
  \item[$\checkmark$] Conda environment \texttt{aiooll}
  \item[$\checkmark$] Local LLMs via Ollama
  \item[$\checkmark$] ML: spam + price prediction
  \item[$\checkmark$] DL: LSTM + MLP on CPU
  \item[$\checkmark$] RAG: local document Q\&A
  \item[$\checkmark$] Agents: LangGraph state machines
  \item[$\checkmark$] CV: face \& motion detection
  \item[$\checkmark$] All tests passing offline
\end{itemize}
\column{0.5\textwidth}
\textbf{Extend AIOOLL yourself:}
\begin{itemize}
  \item Add your own docs to \texttt{src/rag/knowledge\_base/}
  \item Try \texttt{phi3:mini} (3.8B) for better reasoning
  \item Add a new LangGraph agent tool
  \item Plug webcam into the CV motion detector
  \item Submit benchmark results to \texttt{docs/community\_benchmarks.md}
  \item Star the repo and share on LinkedIn!
\end{itemize}
\end{columns}
\vspace{0.5em}
\begin{block}{The AIOOLL Proof}
\centering
\textbf{Intel i3 · 8\,GB RAM · SSD · Ubuntu · Ollama = Complete AI Stack}\\
No GPU. No cloud. No credit card. Just Python and patience.
\end{block}
\end{frame}

%%%%%%%%%%%%%%%%%%%%%%%%%%%%%%%%%%%%%%%%%%%%%%%%%%%%%%%%%%%%%%%%%%%%%%%%%%%%%%%%%%
\begin{frame}[fragile]\frametitle{Quick Reference Card}
\begin{columns}
\column{0.5\textwidth}
\textbf{Environment}
\begin{itemize}\footnotesize
  \item \texttt{conda activate aiooll}
  \item \texttt{ollama serve \&}
  \item \texttt{ollama list}
\end{itemize}
\textbf{Run Modules}
\begin{itemize}\footnotesize
  \item \texttt{make run-ml}
  \item \texttt{make run-dl}
  \item \texttt{make run-rag}
  \item \texttt{make run-agents}
  \item \texttt{make run-cv}
  \item \texttt{make run-llm}
\end{itemize}
\textbf{Launch UIs}
\begin{itemize}\footnotesize
  \item \texttt{make ui-ml}
  \item \texttt{make ui-rag}
  \item \texttt{make ui-agents}
\end{itemize}
\column{0.5\textwidth}
\textbf{Tests}
\begin{itemize}\footnotesize
  \item \texttt{make test}
  \item \texttt{pytest tests/ -v}
  \item \texttt{pytest src/ml/tests/ -v}
\end{itemize}
\textbf{Performance}
\begin{itemize}\footnotesize
  \item \texttt{echo performance | sudo tee /sys/devices/system /cpu/cpu*/cpufreq/scaling\_governor}
  \item \texttt{free -h}
  \item \texttt{btop}
  \item \texttt{sensors}
\end{itemize}
\textbf{Models}
\begin{itemize}\footnotesize
  \item \texttt{ollama pull gemma2:2b}
  \item \texttt{ollama pull qwen2:1.5b}
  \item \texttt{ollama pull nomic-embed-text}
\end{itemize}
\end{columns}
\begin{center}
\large\textbf{github.com/yogeshhk/AIOOLL} 
\end{center}
\end{frame}
